### Conventions

* All graphs are **finite, simple, undirected** (no loops or multiple edges).
* For a graph (X), write (v(X)=|V(X)|) and (e(X)=|E(X)|).
* A “subgraph on (k) vertices” means: choose a subset (S\subseteq V(G)) with (|S|=k), and take **any** subgraph whose vertex set is (S) (equivalently: any graph (J\subseteq G[S])). Then (v(J)=k) and (e(J)\le e(G[S])).

### Ramsey number

For graphs (G,H), the (two–color) Ramsey number (R(G,H)) is the least integer (N) such that every red/blue coloring of the edges of (K_N) contains either:

* a red (not necessarily induced) copy of (G), or
* a blue (not necessarily induced) copy of (H).

Equivalently: for every graph (X) on (N) vertices, either (X) contains (G) as a (not necessarily induced) subgraph, or (\overline{X}) contains (H) as a subgraph.

### The sparsity hypothesis on (G)

The literal condition “every subgraph on (k) vertices has at most (2k-3) edges” is **nonsensical for (k=1)** because it would demand (0\le -1). Standard convention in this area is that such inequalities are intended for (k\ge 2).

So the **minimal corrected statement** is:

> (\forall k\in\mathbb{Z}) with (2\le k\le v(G)), every subgraph (J\subseteq G) with (v(J)=k) satisfies
> [
> e(J)\le 2k-3.
> ]

### “Ramsey size-linear”

“(R(G,H)\ll m)” means: there exists a constant (C_G>0) (depending only on (G)) such that for **every** graph (H) with

* (e(H)=m), and
* no isolated vertices (i.e. every vertex has degree (\ge 1), hence (v(H)\le 2m) and (m\ge 1)),

we have
[
R(G,H)\le C_G,m.
]

### Problem

Assuming the corrected sparsity hypothesis on (G) (for all (k\ge 2)), decide whether:

> (**Claim**) Every such graph (G) is Ramsey size-linear, i.e.
> [
> \exists C_G>0\ \forall H\text{ with no isolates}:\ R(G,H)\le C_G,e(H).
> ]

---
